% Generated draft: requires mathematical and source review.
\documentclass[11pt]{article}
\usepackage[margin=25mm]{geometry}
\usepackage{fontspec}
\setmainfont{Latin Modern Roman}
\usepackage{amsmath,amssymb,mathtools,unicode-math}
\setmathfont{Latin Modern Math}
\usepackage{enumitem,needspace,xurl,hyperref}
\setlength{\parindent}{0pt}
\setlength{\parskip}{5pt}
\setlength{\emergencystretch}{3em}
\newcommand{\N}{\mathbb{N}}
\newcommand{\Z}{\mathbb{Z}}
\newcommand{\Q}{\mathbb{Q}}
\newcommand{\R}{\mathbb{R}}
\newcommand{\C}{\mathbb{C}}
\newcommand{\F}{\mathbb{F}}
\newcommand{\A}{\mathbb{A}}
\newcommand{\PP}{\mathbb{P}}
\newcommand{\E}{\mathbb{E}}
\newcommand{\eps}{\varepsilon}
\newcommand{\dd}{\,\mathrm{d}}
\begin{document}
\section*{UnsolvedMath ID: 20001005}
\textbf{Source problem code:} catalogue entry
\subsection{Definitions and mathematical statement}
Jordan Ellenberg asks for a Deligne--Mumford stack on graphs and, after Matt Baker's comment about a metric-graph paper, asks what harmonic group actions on metric graphs should mean. The source does not specify the category of graphs, families, base topology, stability, or sense of Deligne--Mumford analogue. Baker's paper is not identified.

For the aid's stable moduli interpretation, fix stable weighted marked graph types, with vertex weights \(h(v)\ge0\), prescribed total genus, and stability \(2h(v)-2+\operatorname{val}(v)+\#\text{marks at }v>0\). The reported tropical cone-stack cells are
\[
 [\R_{>0}^{E(G)}/\operatorname{Aut}(G)],
\]
where automorphisms preserve the weights and markings. Cells are glued using edge contractions and appropriate contracted vertex data. This is a cone-stack/tropical analogue with finite combinatorial stabilizers, not an unspecified algebraic Deligne--Mumford stack over schemes. The source's unrestricted graph-stack question is not answered merely by listing the cells without its intended category.

For the metric quotient, let a finite group \(K\) act isometrically on a connected finite metric graph \(\Gamma\) having at least one edge. Subdivide equivariantly to resolve inversions. For an edge \(e\), let \(K_e\) denote its generic point stabilizer, now fixing the edge pointwise; for a point \(x\), let \(K_x\) be its stabilizer. A finite metric-graph map is harmonic at \(x\) if the sum of outgoing integer slopes above a chosen target tangent direction is independent of that target direction; this common sum is its local degree.

The aid reports that giving the orbit edge \(\bar e\) the weighted length
\[
 \ell(\bar e)=|K_e|\ell(e)
\]
makes \(\Gamma\to\Gamma/K\) harmonic, with edge slope \(|K_e|\), local degree \(|K_x|\), and global degree \(|K|\). By contrast, the ordinary path-metric quotient gives slope one on each resolved edge. Its reported criterion is
\[
 \Gamma\longrightarrow\Gamma/K\text{ with the unscaled path metric is harmonic}
 \quad\Longleftrightarrow\quad
 |K_e|=s\text{ is constant over all edges}.
\]
When this holds, its local degree is \(|K_x|/s\) and global degree is \(|K|/s\). The aid says the criterion is subdivision invariant. The distinction is between two quotient metrics on the same orbit combinatorics, not a statement that arbitrary isometric actions fail to admit a harmonic weighted quotient.

The reported obstruction is a \(C_2\)-action on a Y-shaped graph fixing one arm and exchanging two equal-length arms. At the center the unscaled map has one preimage direction over the fixed arm and two over the exchanged-arm orbit, so the directional sums disagree. The generic stabilizers have orders two and one respectively. The weighted metric compensates this difference. The aid calls this a smallest obstruction, but does not state a graph category or minimality measure for that additional claim. These constructions remain supplied partial results requiring verification, not a complete treatment of all graph stacks or harmonic actions.
\subsection{Short English statement}
The source asks about graph moduli stacks and harmonic group actions. The supplied aid describes a stable tropical cone-stack analogue and distinguishes the always available stabilizer-weighted quotient metric from the ordinary orbit path metric, which is harmonic only under a stated constant-stabilizer condition.
\subsection{Sources}
\begin{itemize}
\item Local UnsolvedMath registry, permanent ID \(20001005\), source code \texttt{catalogue entry}; \texttt{aim-combinatorics-notes.json}, \texttt{notes[129]}, tag \texttt{problem}.
\item 2013 AIM workshop ``Generalizations of chip-firing and the critical group,'' Problem 40, Jordan Ellenberg, with Matt Baker's metric-graph comment: \url{https://aimath.org/pastworkshops/chipfiringproblems.pdf}. Statement status \texttt{exact}. No PDF was accessed for this draft.
\item Baker's paper, tropical cone-stack construction, and harmonic-quotient results have no full bibliographic citations in the supplied record. No original remarks or literature field are supplied.
\item Research attempt 1: \texttt{partial\_result}, status \texttt{partially\_solved}, classification \texttt{PARTIAL-PROGRESS}, rigor label \texttt{proved}, novelty confidence low. Created and updated 14 August 2026. This machine-generated aid is not peer reviewed; independently verify before citation.
\item Ulam AI / UnsolvedMath Contributors; CC BY 4.0.
\end{itemize}
\end{document}
